\chapter{Định nghĩa bài toán}

\section{Bối cảnh nghiệp vụ và động lực}

Trong môi trường doanh nghiệp, cuộc họp thường tạo ra quyết định, công việc cần theo dõi và cam kết giữa các thành viên. Tuy nhiên, các thông tin này thường nằm trong audio, ghi chú cá nhân hoặc tin nhắn sau họp. Khi không có hệ thống trích xuất và lưu vết thống nhất, action items dễ bị bỏ sót, deadline không rõ ràng và trách nhiệm có thể bị hiểu sai.

Meeting AI Agent được xây dựng để chuyển audio cuộc họp thành dữ liệu có cấu trúc: transcript theo lượt nói, tóm tắt, action items, người phụ trách, thời hạn, mức ưu tiên và bằng chứng nguồn. Giá trị nghiệp vụ nằm ở việc giảm thời gian viết biên bản họp thủ công, tăng khả năng theo dõi công việc sau họp và giúp người dùng kiểm tra task dựa trên đoạn hội thoại gốc.

\section{Người dùng mục tiêu và các bên liên quan}

Người dùng trực tiếp gồm project manager, technical lead, product owner và thành viên nhóm cần theo dõi công việc sau họp. Các bên liên quan gián tiếp gồm bộ phận vận hành, bảo mật dữ liệu và quản lý chất lượng vì hệ thống xử lý nội dung hội thoại nội bộ, lưu dữ liệu có thể chứa PII và tạo output ảnh hưởng tới workflow quản lý công việc.

\section{Mô tả bài toán và nhu cầu NLP}

Bài toán cần giải quyết không chỉ là speech-to-text. Hệ thống phải nhận diện nội dung công việc trong hội thoại tự nhiên, phân biệt câu thảo luận với cam kết, chuẩn hoá người phụ trách theo roster, diễn giải deadline tương đối theo ngày họp và xuất kết quả thành JSON có thể lưu vào database. Các bước này đòi hỏi NLP ở nhiều mức: ASR, speaker diarization, contextual information extraction, entity resolution và structured generation.

Ví dụ, câu ``Bob, can you send the report by Friday?'' cần được chuyển thành task ``send the report'', assignee được map từ alias ``Bob'' sang worker trong roster, due date được chuẩn hoá theo meeting date, và task giữ source turn để người dùng audit.

\section{Tiêu chí thành công}

Bảng \ref{tab:success-metrics} tóm tắt các chỉ số dùng để đánh giá giá trị nghiệp vụ và chất lượng kỹ thuật của hệ thống.

\begin{table}[H]
\centering
\small
\caption{Tiêu chí thành công về nghiệp vụ và kỹ thuật}
\label{tab:success-metrics}
\begin{tabularx}{\textwidth}{@{}p{3.2cm}X@{}}
\toprule
\textbf{Nhóm chỉ số} & \textbf{Chỉ số} \\
\midrule
Nghiệp vụ & Thời gian tạo minutes sau họp; tỷ lệ task cần sửa sau review; tỷ lệ task có assignee và due date rõ ràng. \\
Kỹ thuật & Precision, recall, F1, assignee accuracy, schema failure rate, hallucination rate, latency trung bình và P95 latency. \\
\bottomrule
\end{tabularx}
\end{table}

Chỉ số nghiệp vụ phản ánh giá trị sử dụng trong quy trình sau họp. Chỉ số kỹ thuật phản ánh chất lượng trích xuất, độ ổn định của đầu ra có cấu trúc và khả năng vận hành của pipeline bất đồng bộ.
